\clearpage\SourceLeaf{367}\NoSourceHead
\HistoricTitle{BRIEFWECHSEL.}
\NumberedItem{1}
\HistoricTitle{\textsc{Gauss} an \textsc{Bessel}. Göttingen 21. October 1810.}
\HistoricCitation{[Briefwechsel zwischen \textsc{Gauss} und \textsc{Bessel}, Leipzig 1880, S. 128--129.]}
\noindent {.}\quad{.}\quad{.}\quad{.}\quad{.}\quad{.}\quad Durch Ihre Untersuchung über das Integral
\[
\int\frac{dx}{\log x}
\]
\noindent haben Sie mir eine grosse Freude gemacht. Ich hatte im Grunde denselben Weg eingeschlagen, bin aber bei weitem nicht so weit darauf vorwärts gegangen, wie Sie\quad Ich setzte nämlich $x=e^y$, wodurch das Integral
\[
\int\frac{e^y\,dy}{y}
\]
\noindent wird. Um diess von $y=a$ bis $y=a+t$ zu integriren, verwandle ich
\[
\frac{e^{a+t}\,dt}{a+t}
\]
\noindent in
\[
\begin{array}{@{}r@{}r@{}r@{}r@{}l@{}}
e^a\,dt\Bigl\{\dfrac{1}{a}&+\dfrac{t}{a}&+\dfrac{tt}{2a}&+\dfrac{t^3}{6a}&+\\[6pt]
&-\dfrac{t}{aa}&-\dfrac{tt}{aa}&-\dfrac{t^3}{2aa}&-\\[6pt]
&&+\dfrac{tt}{a^3}&+\dfrac{t^3}{a^3}&+\text{etc.}\Bigr\}
\end{array}
\]
\noindent also das Integral
\clearpage\SourceLeaf{368}\ArithmeticHead{361}{R}{GAUSS AN OLBERS.}
\[
\begin{aligned}
=e^a\Bigl\{&\frac{1}{a}\left(t+\frac{1}{2}tt+\frac{1}{6}t^3+\text{etc.}\right)\\[6pt]
&-\frac{1}{aa}\left(\frac{1}{2}tt+\frac{1}{3}t^3+\frac{1}{6}t^4+\text{etc.}\right)\\[6pt]
&+\frac{1}{a^3}\left(\frac{1}{3}t^3+\frac{1}{4}t^4+\frac{1}{10}t^5+\text{etc.}\right)\\[4pt]
&\qquad\text{etc.}\qquad\qquad\qquad\Bigr\}
\end{aligned}
\]
\noindent wo man die Coefficienten von $\dfrac{1}{a},\ \dfrac{1}{aa},\ \dfrac{1}{a^3}$ etc. für bestimmte Werthe von $t$ ein für allemal berechnen kann. Offenbar kann man auch $t$ negativ setzen, und dann kommt die Operation mit der Ihrigen überein. Nur hatte ich mich beschränkt, für $t$ rationale Werthe $\frac{1}{2}$ und $1$ zu substituiren, und so die Integrale für $y=5,\ 6,\ 7,\ 8,\ 9,\ 10,\ 11$ bestimmt\quad {.}\quad{.}\quad{.}\quad{.}\quad{.}\quad{.}\quad{.}\quad{.}\quad{.}\quad{.}\quad{.}\quad{.}
\par
